\markboth{STATE OF THE ART}{STATE OF THE ART}
\section{State of the Art}\label{sec:sota}

This section reviews the four research pillars Training Copilot builds upon: agentic AI and multi-agent systems, tool-integration protocols, sports analytics and wearable technology, and retrieval-augmented generation.

\subsection{Agentic AI and Multi-Agent Systems}\label{sec:sota:agents}

Large language models have grown quickly, shifting the field from static text generation toward \textit{agentic} systems that plan, reason, and act in external environments. The Transformer architecture~\citep{vaswani_attention_2017}, chain-of-thought prompting~\citep{wei_chain-of-thought_2022}, and the scaling of models to hundreds of billions of parameters~\citep{brown_gpt3_2020} showed that LLMs can reason step by step once prompted to write out intermediate steps, and later frontier models reached human-level performance on professional and academic benchmarks~\citep{openai_gpt4_2023}. We organise the literature along three themes: reasoning architectures, agent-internal structure, and multi-agent coordination.

\subsubsection{Reasoning Architectures}

Two complementary patterns ground LLM reasoning in external action. \cite{yao_react_2023} introduced \textit{ReAct}, which interleaves reasoning traces (``thought'' steps) with task-specific actions in a single loop. The model generates a thought, executes an action (for example \texttt{search[entity]} against a Wikipedia API), observes the result, and repeats until done. Because every claim is grounded in a retrieved observation rather than generated from parameters, ReAct sharply cut hallucinated reasoning compared to chain-of-thought prompting. Training Copilot's specialists apply this pattern directly, as each is implemented as a LangGraph ReAct agent whose ``actions'' are MCP tool calls. \cite{shinn_reflexion_2023} extended the loop across episodes through \textit{Reflexion}. After a failed attempt, the agent records a short note in plain language and retains it for the next attempt, so it improves without retraining. We use this idea in one subsystem, where the chart generator retries once with the error message as feedback (Section~\ref{sec:architecture:supporting}). However, both patterns assume a single agent operating over a fixed tool set and address neither multi-agent coordination nor dynamic tool discovery across different sources.

\subsubsection{Agent-Internal Architecture}

\cite{park_generative_2023} showed that LLM agents behave believably when given a \textit{memory stream}, a \textit{reflection} step, and a \textit{planning} step. \cite{wang_survey_2024} group these ideas into four parts that recur across agent designs: profiling, memory, planning, and action. Both works consider agents in isolation but leave the coordination of agents with distinct roles and tool sets unaddressed.

\subsubsection{Multi-Agent Coordination}

Several works illustrate the multi-agent design space. \cite{shen_hugginggpt_2023} proposed HuggingGPT, where a central controller decomposes a request, selects specialist models, executes them, and summarises the results; this is an orchestrator-specialist decomposition, applied to model selection rather than tool-based data retrieval. \cite{wu_autogen_2023} introduced AutoGen and found that several role-specialised agents with tool access beat a single monolithic agent on complex tasks, and that different conversation patterns (sequential, parallel, or nested) give different ways to coordinate them. \cite{guo_large_2024} survey recurring open challenges in orchestration, inter-agent communication, and scalability, and \cite{acharya_agentic_2025} distinguish agentic AI by its capacity to make autonomous, goal-directed decisions and adapt as conditions change. Across these works, domain-scoped specialists outperform monolithic agents on mixed tasks; none of them, however, use standardised tool discovery or inter-agent protocols, and each builds its own ad-hoc registry and communication layer.

\subsection{Tool Integration Protocols}\label{sec:sota:protocols}

LLM agents are only useful if they can reach external tools and data. The work includes the three stages: teaching models to use tools at all, standardising how tools are discovered and invoked, and standardising how tool-using agents talk to each other.

\subsubsection{Tool-Augmented Language Models}

\cite{schick_toolformer_2023} showed that a model can learn to use tools on its own. It inserts candidate API calls into ordinary training text, runs them, and keeps a call only when its result lowers the model's loss on the next tokens. The trained model chooses tools well, but it can only make one call per example and cannot rephrase a query when the result is unhelpful. \cite{patil_gorilla_2023} studied tool selection at a larger scale. They fine-tuned a retrieval-aware model on 1{,}645 APIs from three model hubs, and it called far fewer non-existent APIs than a much larger general-purpose model did. \cite{qu_tool_2024} split tool use into four stages (planning, selection, calling, and response generation) and show, on benchmarks with over 16{,}000 tools, that once there are more than a few hundred tools you can no longer fit all their descriptions into the prompt, so a retrieval step has to shortlist the relevant ones first. We use this directly: each specialist sees at most three MCP servers (1-41 tools), which keeps its tool set well under that limit.

\subsubsection{From Ad-hoc Integration to Standardised Protocols}

The proliferation of tool-augmented agents created an $M \times N$ fragmentation problem, where $M$ applications each integrate $N$ data sources separately. Two complementary protocols address this. The \textbf{Model Context Protocol (MCP)}~\citep{anthropic_mcp_2024, anthropic_mcp_announcement_2024} defines a universal client-server protocol over JSON-RPC\,2.0~\citep{jsonrpc_spec_2013} for tool discovery and invocation. It exposes Prompts, Resources, and Tools, and separates authentication from tool declaration, so a host can use arbitrary servers without provider-specific code. \cite{jayanti_enhancing_2026} proposed Context-Aware MCP (CA-MCP), which adds a \textit{Shared Context Store} that servers read and write after the central LLM seeds it with goals and constraints, cutting execution time and orchestration calls on a planning benchmark. Training Copilot does not adopt this shared-memory pattern (A2A handles inter-agent coordination instead), but the result supports MCP's extensibility. The \textbf{Agent-to-Agent protocol (A2A)}~\citep{google_a2a_2025}, donated to the Linux Foundation after Google's April 2025 announcement, standardises inter-agent communication over JSON-RPC\,2.0 around \textit{opaque internals}. Agents advertise their capabilities through Agent Cards but expose only results, not their reasoning, which enables composition across frameworks and providers. \cite{yang_survey_2025} distinguish the two as complementary: MCP is \textit{context-oriented} (agents to data), and A2A is \textit{coordination-oriented} (agents to agents). Both protocols are new enough that their security surface is still being charted. \cite{hou_mcp_2025} give the first systematic analysis of MCP's threat model across the full server lifecycle (creation, deployment, operation, and maintenance), identifying tool poisoning and injection risks that any multi-tenant deployment must address; we scope these safeguards out of this prototype and return to them in Section~\ref{sec:discussion:limitations}. We are not aware of a published system that combines both protocols in a domain-specific multi-agent application.

\subsection{Sports Analytics and Wearable Technology}\label{sec:sota:sports}

Sports science agrees that good training management means combining several physiological and contextual signals, but no consumer tool actually reasons across them. We look at three themes: training load and periodisation, recovery monitoring, and wearable reliability.

\subsubsection{Training Load: Monitoring and Injury Prevention}

\cite{bourdon_monitoring_2017}, in an international consensus statement (11 authors, five countries), distinguish \textit{external load} (objective work, measured by GPS and accelerometers) from \textit{internal load} (the body's response: heart rate, HRV, session RPE, lactate), and recommend monitoring both together, since the same external load produces different internal responses depending on fatigue or illness. This ``uncoupling'' is itself diagnostic. They establish the acute-to-chronic workload ratio (ACWR), with 0.8-1.3 as low-risk, and \cite{halson_monitoring_2014} add that no single biomarker is sufficiently validated on its own. \cite{gabbett_training_2016} formalise the ``training-injury prevention paradox'' from cricket, Australian football, and rugby league data: well-conditioned athletes have \textit{fewer} injuries than undertrained ones, since conditioning itself protects, while injury risk rises sharply when weekly load spikes above the chronic baseline instead of staying roughly constant. \cite{grgic_periodization_2017} extend this to periodisation, finding in a meta-analysis that linear and undulating schemes produce equivalent hypertrophy, so the specific periodisation pattern does not significantly affect outcomes. Together these findings show that load management needs trend analysis over time, not point-in-time metrics. Training Copilot's load specialist computes chronic and acute training load (CTL, ATL) and their difference, training-stress balance (TSB), from Strava data (derived from the fitness-fatigue impulse-response model of \cite{morton_modeling_1990}) to put this into practice. While the ACWR itself has faced methodological criticism, \cite{impellizzeri_training_2020} identified mathematical coupling (acute load appears in both numerator and denominator) and showed that simulations can produce spurious ACWR-injury associations from unrelated data. This means, our system presents the ACWR as one contextual signal rather than a standalone predictor.

\subsubsection{Recovery: From Single Metrics to Multi-Signal Assessment}

\cite{kellmann_recovery_2018} argue in a consensus statement that systematic recovery monitoring (not just load tracking) should be integral to training, since the stress-recovery balance determines long-term sustainability. HRV has become a central biomarker, but its interpretation is not straightforward. \cite{plews_training_2013}, using data from Olympic and World Champion rowers, show that in elite athletes both increases \textit{and} decreases in HRV can signal negative adaptation, because vagal HRV saturates at very low resting heart rates. They recommend weekly rolling averages of Ln\,rMSSD over single-day values, interpreted against each athlete's own baseline; individual profiles vary so widely that a fixed-threshold system would misread them. \cite{lundstrom_hrv_2023} review HRV-guided training in consumer settings and note the gap between clinical- and consumer-grade measurement. The strongest evidence for multi-signal approaches comes from \cite{rothschild_predicting_2024}, who tracked 43 endurance athletes over 12 weeks and tested nine machine-learning models predicting daily perceived recovery. Prediction accuracy was only moderate and varied widely across individuals, while soreness, sleep, and prior-day status proved more predictive than HRV alone. \cite{wackerhage_personalized_2021} frame this difficulty as inherent, since a training plan is a complex mix of interacting, time-varying interventions that makes purely evidence-based decisions impractical. This motivates tool-assisted approaches like Training Copilot's recovery specialist, which queries HRV, sleep, Body Battery, and stress together to approximate this multi-signal assessment.

\subsubsection{Wearable Reliability and Data Ecosystem Biases}

Wearable data reliability is a prerequisite for any analytics system built on it. \cite{fuller_reliability_2020}, reviewing 158 publications across 45 device models from nine manufacturers, found step counts and heart rate reasonably accurate at rest but degrading during high-intensity exercise and in free-living settings, with energy expenditure inaccurate across all devices and heart-rate accuracy varying markedly by brand. \cite{seckin_review_2023} note that translating raw sensor data into actionable coaching remains an open problem, and \cite{vanhooren_wearables_2020} identify challenges in feedback timing and injury-risk communication for real-time wearable feedback in running. Broader reviews of wearable performance devices in sports medicine~\citep{li_wearable_2016} and of machine-learning-driven biomechanical analytics~\citep{alzahrani_advanced_2024} reach a similar conclusion: sensing has outpaced interpretation, and the real bottleneck is turning multi-stream data into decisions. \cite{venter_bias_2023} further quantify selection bias in Strava's crowdsourced data (which overrepresents recreationally active, male, and middle-aged users), a bias professional sports avoid~\citep{bourdon_monitoring_2017} but consumer platforms must acknowledge. Training Copilot mitigates this partially through multi-source corroboration (cross-referencing HRV, sleep, and Body Battery), but cannot fully compensate for sensor inaccuracy. Across these themes the gap is in synthesis, not availability: the tools that \textit{collect} this data do not \textit{reason across} it.

\subsection{Retrieval-Augmented Generation}\label{sec:sota:rag}

Retrieval-Augmented Generation pairs a retriever with a generator, so answers are grounded in retrieved text rather than only in the model's parameters. \cite{lewis_retrieval-augmented_2020} introduced the original architecture (a generator augmented with a Dense Passage Retriever, DPR, over a FAISS index of Wikipedia passages) in two variants: RAG-Sequence, where one retrieved passage conditions the whole output, and RAG-Token, where a different passage can condition each token. RAG set the state of the art on open-domain QA and was judged markedly more factual than a generator without retrieval, and its retrieval index can be swapped without retraining. The retrieval component builds on \cite{karpukhin_dense_2020}, whose dual-encoder DPR clearly outperforms sparse keyword retrieval, and on \cite{reimers_sentence-bert_2019}'s Sentence-BERT, which adapts BERT with a Siamese objective for efficient cosine-similarity comparison. The \texttt{sentence-transformers} library built on this work provides Training Copilot's embedding model (\texttt{all-MiniLM-L6-v2}, $\sim$90\,MB, running locally without an embedding API). \cite{gao_retrieval-augmented_2024} survey the progression from Na\"ive RAG (retrieve then generate) through Advanced RAG (query rewriting, re-ranking) to Modular RAG (pluggable components). Training Copilot implements the na\"ive form (embed, retrieve by cosine similarity, generate with the retrieved context), which suffices for its curated corpus of $\sim$11{,}000 chunks from five contemporary German-language sports-science textbooks.

\subsection{AI-Assisted Sports Coaching}\label{sec:sota:coaching}

AI-assisted sports coaching is an emerging field where three limitations recur: insufficient access to the athlete's actual data, limited cross-domain reasoning, and single-agent architectures that do not scale to heterogeneous tool sets.

Regarding the data access and personalisation, \cite{monteiro-guerra_personalization_2020} reviewed 28 papers on real-time physical-activity coaching apps across seven personalisation dimensions. Feedback and goal setting were nearly universal, whereas context-awareness and behavioural adaptation appeared in only a handful of systems. \cite{luo_promoting_2021} reach a similar conclusion for conversational agents specifically (chatbots sustain engagement, but most lack wearable-data access), and \cite{puce_harnessing_2025} confirm this in a systematic review of generative AI for exercise prescription. Although AI-generated plans followed foundational training principles, they lacked specificity, progression, and, most importantly, real-time physiological feedback. Identical prompts also sometimes produced different outputs.

On architecture, the most directly relevant prior system is GPTCoach~\citep{joerke_gptcoach_2024}, an LLM-based coaching chatbot evaluated at CHI~2025. It pairs GPT-4 with a prompt-chaining design (separate calls for dialogue-state classification, motivational-interviewing strategy, and tool-call decisions) and two tools querying three months of Apple HealthKit data via HL7 FHIR. A small lab study rated it highly for feeling supported and for empathy, and external coding found almost all of its behaviour motivational-interviewing-consistent or neutral. GPTCoach is still a single agent with two tools and one data source, however; it integrates no weather, calendar, or route data, and was evaluated only on a single onboarding session. At the architectural level, \cite{vemuri_multiagent_2021} proposed SMART-ALEC, a multi-agent health-coaching system with a per-user analyser and communicator module, prototyped as an SMS system with just-in-time interventions on Apple Watch. It predates MCP and A2A and coordinates through ad-hoc commitment exchange, but it nonetheless supports the modular-specialist-with-central-coordinator philosophy that Training Copilot extends with standardised protocols and LLM reasoning in place of finite-state dialogues.

On the human-AI boundary, \cite{barger_artificial_2025} ran a mixed-methods RCT comparing sessions where participants believed they were coached by an AI (actually a human behind a live-motion avatar) against sessions coached openly by a human: working-alliance scores did not differ significantly, and participants in the AI condition used fewer impression-management behaviours and disclosed more sensitive topics. This supports a synergy model of narrowly focused AI coaches working alongside humans. \cite{gabarron_human_2024} survey AI-driven solutions for physical activity more broadly, finding recommender systems the most studied, followed by conversational agents, with sparse long-term evidence for either.

\subsection{Summary and Research Gap}\label{sec:sota:summary}

Each of these research streams has matured on its own, but they share one gap: none brings together the other capabilities that integrated training support needs. Multi-agent AI often relies on ad-hoc registries rather than standardised protocols. MCP and A2A solve the $M \times N$ integration problem in principle but have not been combined in a domain-specific application. Sports science shows that training decisions require the synthesis of load, recovery, and contextual data, whereas consumer platforms typically present these signals in isolation. AI-coaching systems therefore continue to trade data access against reasoning capability. RAG grounds answers in external sources but has not yet been combined with live physiological data for coaching.

Table~\ref{tab:sota_comparison} in Appendix~\ref{sec:appendix:comparisons} makes this concrete, scoring the most directly comparable prior systems (academic multi-agent frameworks, the closest published MCP extension, and the leading commercial AI coaching products) against four capabilities: multi-agent coordination, standardised protocols, multi-source live data, and retrieval-grounded knowledge, plus the application domain each targets. Its key finding: no prior system combines all four capabilities, and none of the few sports- or health-domain systems demonstrates more than one of them.

The research gap is therefore not in any individual capability but in their integration: no existing system combines standardised multi-source access (MCP), multi-agent coordination (A2A), domain-scoped specialist reasoning (LangGraph ReAct), and retrieval-grounded knowledge (RAG) for athlete-facing training support. Training Copilot's architecture maps each of these limitations to a specific design decision, detailed in Section~\ref{sec:architecture}.
